\chapter{نظام الملفات}
\label{CH:FS}

تتمثل الغرض من نظام الملفات في تنظيم وتخزين البيانات على وسائل التخزين غير المتطايرة (مثل القرص الصلب أو SSD)، ودعم مشاركة البيانات بين البرامج، وتوفير القدرة على استرجاع الملفات بأسماء معبرة.

يقدم هذا الفصل الطبقات السبع المنظمة لنظام الملفات في xv6، وطريقة عمل ذاكرة التخزين المؤقت المخبئية (Buffer Cache)، وحصص كتل القرص، وعقد الفهرسة (Inodes)، والدلائل، وواصفات الملفات.

\section{نظرة عامة وطبقات نظام الملفات}

ينقسم نظام الملفات في xv6 إلى سبع طبقات متسلسلة:
\begin{enumerate}
\item \textbf{القرص (Disk)}: الكتل العتادية المادية المكونة للتخزين.
\item \textbf{المخزن المؤقت (Buffer Cache)}: تزامن الوصول لكتل القرص وتخزينها في الذاكرة لتسريع القراءة والتأكد من تحريرها.
\item \textbf{السجل الدوري (Logging)}: ضمان تنفيذ المعاملات المركبة ككتلة ذرية للتعافي عند انقطاع التيار الكهربائي.
\item \textbf{عقد الفهرسة (Inode Layer)}: توفير ملفات بدون أسماء ومتابعة مواضع الكتل والميتاداتا.
\item \textbf{الدلائل (Directory Layer)}: ربط الأسماء النصية برقم عقدة الفهرسة (Inode Numbers).
\item \textbf{المسارات (Path Names)}: فك وتتبع المسارات الشجرية الممتدة (مثل \lstinline{/usr/bin/sh}).
\item \textbf{واصفات الملفات (File Descriptors)}: تجريد كائنات النظام (الملفات، الأنابيب، الأجهزة) لبرامج المستخدم.
\end{enumerate}

يُقسم القرص المادي في xv6 إلى كتل بحجم 1024 بايت (Block 0 للإقلاع، Block 1 ككتلة سوبر Superblock تحتوي على حجم ونظام الملفات، يليها السجل الدوري Log، ثم عقد الفهرسة Inodes، ثم خريطة البتات للكتل الحرة Bitmap، وأخيراً كتل البيانات Data Blocks).

\section{طبقة التخزين المؤقت المخبئي}

تحقق طبقة المخزن المؤقت (\indextext{Buffer Cache}) هدفين رئيسين:
\begin1. ميكانيكية التزامن: ضمان وجود نسخة واحدة فقط لكل كتلة قرص في الذاكرة، وأن نواتًا واحدة فقط تعدل محتوى الكتلة في وقت محدد.
2. الكفاءة والتسريع: التخزين المخبئي للكتل الأكثر استخدامًا في الذاكرة لتجنب القراءة البطيئة المكررة من القرص.

\section{الكود البرمجي: التخزين المؤقت المخبئي}

تُدار الكتل عبر القائمة الدائرية المزدوجة \lstinline{bcache} في \lstinline{kernel/bio.c}:
\begin{itemize}
\item \lstinline{bread(dev, blockno)}: تجلب الكتلة المطلوب قراءتها من المخزن المؤقت، أو تقرأها من القرص عبر محرك virtio وتثبيتها.
\item \lstinline{bwrite(b)}: تكتب البيانات المعدلة في الكتلة فورًا إلى القرص المادي.
\item \lstinline{brelse(b)}: تحرر الكتلة وتضعها في أصل القائمة وفق آلية LRU (الأقل استخدامًا مؤخرًا - Least Recently Used).
\end{itemize}

\section{الكود البرمجي: مخصص كتل القرص}

تُتابع الكتل المتاحة للبيانات عبر خريطة البتات (Bitmap Blocks).
تضم الدالتان التاليان في \lstinline{kernel/fs.c} إدارة الكتل:
\begin{itemize}
\item \lstinline{balloc(dev)}: تبحث عن بت محايد \lstinline{0} في خريطة البتات، وتغيره إلى \lstinline{1} وتصفّر الكتلة المخصصة وتكتُب التغيير على القرص.
\item \lstinline{bfree(dev, b)}: تعيد البت المخصص ليكون \lstinline{0} في خريطة البتات.
\end{itemize}

\section{طبقة عقد الفهرسة (Inode Layer)}

مصطلح \indextext{عقدة الفهرسة} (Inode) يُعبر عن الهيكل الحافظ لبيانات الملف الميتا (حجم الملف، نوعه، عدد الوصلات الصلبة) وقائمة كتل البيانات المحتوية على محتواه.

تحتوي عقدة الفهرسة في xv6 على:
\begin{itemize}
\item 12 عنوانًا مباشرًا (Direct Blocks): تشير مباشرة لـ 12 كتلة بيانات ($12 \times 1024 = 12$ كيلوبايت).
\item عنوان غير مباشر واحد (Indirect Block): يشير لكتلة تحتوي على 256 عنوانًا لكتل بيانات إضافية ($256 \times 1024 = 256$ كيلوبايت).
\item الحد الأقصى لحجم الملف في xv6 هو $12 + 256 = 268$ كتلة ($268$ كيلوبايت).
\end{itemize}

\section{الكود البرمجي: Inodes}

تُحفظ Inodes ماديًا على القرص، وتُنسخ في ذاكرة النواة عبر الهيكل \lstinline{struct inode}:
\begin{itemize}
\item \lstinline{ialloc(dev, type)}: تخصص Inode جديدة على القرص.
\item \lstinline{iget(dev, inum)}: تجلب Inode من الذاكرة أو تقرأها من القرص وتزيد معرّف التثبيت \lstinline{ref}.
\item \lstinline{ilock(ip)}: تحوز قفل النوم الخاص بـ Inode وتضمن قراءة بياناتها الحالية من القرص.
\item \lstinline{iunlock(ip)}: تحرر القفل الخاص بـ Inode.
\item \lstinline{iput(ip)}: تنقص عدد التثبيتات \lstinline{ref}. وإذا وصل إلى الصفر دون وجود أي وصلات صلبة (\lstinline{nlink == 0})، تحذف Inode وتحرر كتل بياناتها عبر \lstinline{itrunc()}.
\end{itemize}

\section{الكود البرمجي: محتوى عقد الفهرسة}

تنفذ الدالتان \lstinline{readi()} و \lstinline{writei()} القراءة والكتابة داخل كتل Inode عبر تحويل إزاحة البايت إلى رقم الكتلة المقابل (باستخدام \lstinline{bmap()}) ونقل البيانات من وإلى الذاكرة المؤقتة.

\section{الكود البرمجي: طبقة الدلائل}

الدليل هو Inode خاصة يحتوي محتواها على مصفوفة من المدخلات الدليلية \lstinline{struct dirent}.
تحتوي كل مدخلة على:
\begin{itemize}
\item رقم Inode (\lstinline{ushort inum}).
\item اسم الملف النصي (\lstinline{char name[DIRSIZ]}، حيث DIRSIZ = 14 حرفًا).
\end{itemize}
تُستخدم الدالة \lstinline{dirlookup()} للبحث عن اسم ملف داخل الدليل وإعادة Inode الخاصة به.

\section{الكود البرمجي: مسارات الملفات}

يقوم التابع \lstinline{namei(path)} بفك المسارات النصية (مثل \lstinline{/a/b/c}) عن طريق البدء من الدليل الجذري \lstinline{root} (إذا بدأ المسار بـ \lstinline{/}) أو الدليل الحالي \lstinline{cwd}، والاستدعاء المتكرر لـ \lstinline{dirlookup()} لكل جزء من مكونات المسار حتى الوصول للملف المستهدف.

\section{طبقة واصفات الملفات}

تُربط الملفات المفتوحة بالعمليات من خلال جدول واصفات الملفات في العملية (\lstinline{p->ofile[]}) والجدول العام للملفات المفتوحة في النواة (\lstinline{ftable}).
يحتوي الهيكل \lstinline{struct file} على نوع الكائن (ملف عادي، أنبوب، أو جهاز)، ومؤشر التزحيق الحالي (\lstinline{off})، وInod المرفق.

\section{الكود البرمجي: استدعاءات النظام}

تنفذ البرمجيات في \lstinline{kernel/sysfile.c} استدعاءات النظام الخاصة بنظام الملفات (\lstinline{sys_open}، \lstinline{sys_read}، \lstinline{sys_write}، \lstinline{sys_link}، \lstinline{sys_unlink}، \lstinline{sys_mkdir}) وتجمع التعديلات داخل معاملات السجل الدوري (\lstinline{begin_op()} و \lstinline{end_op()}).

\section{العالم الحقيقي}

تستخدم أنظمة الملفات الإنتاجية الحديثة (مثل ext4 و ZFS و Btrfs) تقنيات متطورة للغاية مثل الأشجار المتوازنة (B-Trees) لتسريع البحث في الدلائل الضخمة، والنسخ عند الكتابة (Copy-On-Write Snapshots)، وتخصيص المدى المتصل (Extents) بدلاً من الكتل المنفصلة لتسريع الأداء وتفادي التجزئة.

\section{تمارين}

1. ما الحد الأقصى لحجم الملف في xv6 إذا أضفنا عنوانًا غير مباشر مزدوجًا (Double Indirect Block)؟
2. أضف دعمًا للملفات ذات الأسماء الطويلة التي تتجاوز 14 حرفًا في xv6.
